% Generated draft: requires mathematical and source review.
\documentclass[11pt]{article}
\usepackage[margin=25mm]{geometry}
\usepackage{fontspec}
\setmainfont{Latin Modern Roman}
\usepackage{amsmath,amssymb,mathtools,unicode-math}
\setmathfont{Latin Modern Math}
\usepackage{enumitem,needspace,xurl,hyperref}
\setlength{\parindent}{0pt}
\setlength{\parskip}{5pt}
\setlength{\emergencystretch}{3em}
\newcommand{\N}{\mathbb{N}}
\newcommand{\Z}{\mathbb{Z}}
\newcommand{\Q}{\mathbb{Q}}
\newcommand{\R}{\mathbb{R}}
\newcommand{\C}{\mathbb{C}}
\newcommand{\F}{\mathbb{F}}
\newcommand{\A}{\mathbb{A}}
\newcommand{\PP}{\mathbb{P}}
\newcommand{\E}{\mathbb{E}}
\newcommand{\eps}{\varepsilon}
\newcommand{\dd}{\,\mathrm{d}}
\begin{document}
\section*{UnsolvedMath ID: 700300}
\textbf{Source problem code:} catalogue entry
\subsection{Definitions and mathematical statement}
For a field \(K\) and a positive integer \(n\), let \(\mathrm{GL}_n(K)\) denote the group of invertible \(n\times n\) matrices over \(K\). A group is linear if it admits an injective homomorphism into some \(\mathrm{GL}_n(K)\). Under the usual interpretation, a group \(G\) is locally linear if every finitely generated subgroup \(H\leq G\) is linear; the field and dimension may depend on \(H\). The supplied record does not explicitly specify this convention or impose a common field or characteristic, so this interpretation requires confirmation against the source.

A group is countable if its underlying set is finite or countably infinite. A group \(P\) is finitely presented if \(P\cong F(X)/\langle\!\langle R\rangle\!\rangle\), where \(X\) is a finite set, \(F(X)\) is the free group on \(X\), and \(R\subseteq F(X)\) is finite.

\textbf{Question (F. B. Cannonito and C. F. Miller III).} Under the interpretation above, is it true that for every countable locally linear group \(G\), there exist a finitely presented group \(P\) and an injective homomorphism \(\iota:G\hookrightarrow P\)?

\textbf{Supplied background result.} The record attributes to G. Baumslag, F. B. Cannonito, and C. F. Miller III, \emph{Math. Z.} \textbf{153} (1977), 117--134, the result that every countable group which is locally linear of bounded degree embeds in a finitely presented group with solvable word problem. On the corresponding usual reading, bounded degree means that there exists a positive integer \(d\) such that every finitely generated subgroup \(H\leq G\) embeds in \(\mathrm{GL}_{n_H}(K_H)\) for some field \(K_H\) and some \(1\leq n_H\leq d\). The record does not specify the precise field convention for this background result either. A finitely presented group has solvable word problem if an algorithm decides, for every word in its finite generators and their inverses, whether that word represents the identity. The question itself imposes neither a uniform degree bound on \(G\) nor a solvable-word-problem requirement on \(P\).
\subsection{Short English statement}
Does every countable group whose finitely generated subgroups are linear embed in a finitely presented group? The supplied record reports an affirmative result under a uniform bound on the matrix dimensions, with an ambient finitely presented group having solvable word problem. The precise field conventions in both formulations require source confirmation.
\subsection{Sources}
\begin{itemize}
\item Local UnsolvedMath registry, supplied complete source record: permanent ID \(700300\), exact source problem code \texttt{catalogue entry}, title ``Kourovka Notebook Problem 5.16,'' upstream revision \texttt{d371ee917ac9647fbb06cd6fc529deec6f3dda82}. The registry labels the problem \texttt{open}; this is a supplied status, not an independently verified literature assessment.
\item \emph{The Kourovka Notebook (Unsolved Problems in Group Theory)}, Problem 5.16, proposed by F. B. Cannonito and C. F. Miller III in the 5th issue (1976). The supplied record identifies the 21st issue edition (2026), \texttt{arXiv:1401.0300v48}; its collection description names E. I. Khukhro and V. D. Mazurov as editors.
\item G. Baumslag, F. B. Cannonito, and C. F. Miller III, \emph{Math. Z.} \textbf{153} (1977), 117--134. This is the supplied citation for the bounded-degree embedding result; no article title is supplied.
\end{itemize}
\end{document}
